%%%PAGE 187 rot=0 legibility=fair summary=粘贴印刷纸（盖住笔记页中段）：第20题改装电压表并测电池电动势和内阻（图甲乙丙）及手写批注；纸外上沿/下沿笔迹与p186为同一笔记页，未重复转录
% 说明：本页又是 p186 同一笔记页的另一张照片，中段被另一张粘贴的印刷纸（紫色，与 p185 盖住的不是同一张）盖住，因此：
%  (1) 纸外上沿（验证力的平行四边形定则笔记、钩码链图、行间“大小方向”及右侧淡字）与下沿（“直流”装置图残部、利用第n点与起始点计算、v^2=2gh、(1/2)v^2=gh、图线斜率表示重力加速度g）的手写内容，
%      均与 p186 相同，已在 p186.tex 完整转录，此处不重复；
%  (2) 本页独有的新内容是粘贴印刷纸上的第20题及其手写圈画/批注，转录如下。汇编时请勿把 p187 上沿/下沿与 p186 重复收录。
%%%BLOCK id=187-01 ch=10 sec=电表改装·改装电压表并测电源电动势内阻 kind=problem alt=none cont=none y=0.165-0.80 exp=1
% 粘贴的印刷题（紫色纸）；原稿圈画“改装后的”，划线“内阻R_V约为5kΩ”“某时刻电压表的示数为80V”“保持滑片P的位置不变，调节电阻箱，使电压表的示数为”“其测量值总是”“示数记为U）连接成如图乙所”“电阻箱接入电路的阻值R和电压表的示数U”
\begin{problem}[粘贴印刷题 20]
20．电动汽车充电比燃烧汽油更便宜，受到市场的欢迎。某兴趣小组为了测量电动汽车上电池的电动势 $E$（$300\unit{V}\sim400\unit{V}$）和内阻 $r$（$0\sim10\,\Omega$），需要将一个量程为 $100\unit{V}$ 的电压表（内阻 $R_V$ 约为 $5\unit{k\Omega}$）改装成量程为 $400\unit{V}$ 的电压表，然后再测量电池的电动势和内阻。以下是该实验的操作过程。

\noindent（题干上方行间有淡铅笔小字草算，字迹模糊，可见 \unclear{$\frac12(n\cdots)^2$} 等）

%FIG p187_a 0.012 0.265 0.287 0.46
\notefig[0.3]{figures/p187_a.png}

%FIG p187_b 0.288 0.335 0.53 0.463
\notefig[0.28]{figures/p187_b.png}

%FIG p187_c 0.54 0.31 0.81 0.445
\notefig[0.28]{figures/p187_c.png}

（图丙：$\dfrac1R$--$\dfrac1U$ 图线，图中标注见图）

\noindent\textbf{图旁手写：}图甲上：$20\unit{V}$（V 表上方）、$60\unit{V}$（$R_p$ 上方）、$80\unit{V}$（V 表左上方与下方导线旁各一处，字迹反复涂描）；V 与 $R_p$ 被一个椭圆圈在一起，其右写 $R_x$。图乙上：电压表圈内手写 $U$，$R_p$ 右侧手写 $R$。

串联分压、并联分流、\unclear{欧姆}定律% 位于图乙/图丙上方

$r_V=3.5\unit{k\Omega}$% CHECK: 3 与 5 之间可能漏小数点，原稿写作“3 5kΩ”
\quad 其上方另有 \unclear{5…}、\unclear{$\triangle$…$\le$}（字迹潦草）

分压法选小阻\quad $U$ 与滑片位置成线性关系

$\dfrac{1.8\times10^{-8}}{21}$\quad \unclear{$2.84\times10^{-9}$}\quad（页面右缘另有草算 \unclear{22、484、21…}）

（1）由于不知道该电压表内阻的确切值，该兴趣小组将一个最大阻值为 $50\unit{k\Omega}$ 的电阻箱 $R_p$ 与电压表串联后，利用如图甲所示的电路进行改装。将滑动变阻器的滑片 $P$ 移至右端，同时将电阻箱的阻值调为零。闭合开关 S，将滑片 $P$ 向左移动，某时刻电压表的示数为 $80\unit{V}$。保持滑片 $P$ 的位置不变，调节电阻箱，使电压表的示数为 \underline{\hspace{2.5em}} $\unit{V}$．……再改变电阻箱的阻值，保持电压表和电阻箱串联，撤去其他电路就得到改装后的电压表。% 此行右端“V．”之后被页边截断

\noindent\textbf{（1）问处手写：}“$80\unit{V}$”上方行间：分压\unclear{与 $R$ 有关，与 $R_x$ 无关}；“保持滑片 $P$ 的位置不变”下方：分压不变；闭合开关 S 上方：\unclear{7…}；空格处手写 \unclear{20}，被圈并叉掉，其右上另有 \unclear{$\frac1{20}$}。

（2）用改装后的电压表接入电路测量已知电压时，由于系统误差，其测量值总是 \struck{大于}（填“大于”、“等于”或“小于”）真实值。% 原稿：“大于”为手写，被圈出并斜线划去

\noindent\textbf{（2）问处手写：}小于（写在“大于”右上方）；实际大于 $400\unit{V}$；页边右侧另有圈出的 \unclear{阻}。

（3）该兴趣小组利用另一个电阻箱 $R$ 和改装后的电压表（电压表的表盘没有改变，示数记为 $U$）连接成如图乙所示的电路来测量该电池的电动势和内阻。根据测得的电阻箱接入电路的阻值 $R$ 和电压表的示数 $U$ 作出 $\dfrac1R-\dfrac1U$ 图像如图丙所示，则该电池的电动势 $E=$ \underline{\;36\unclear{0}\;} V、内阻 $r=$ \underline{\;9\;} $\Omega$。% CHECK: E 的空格内手写“360”，末位被粘贴纸边缘遮挡

\noindent 页面下沿贴着纸边处另有手写 \unclear{$4U$…}（后面被纸边遮挡）。
\end{problem}
%%%END
%%%PAGEEND 187
